\documentclass[11pt,a4paper]{article}

\usepackage[utf8]{inputenc}
\usepackage[T1]{fontenc}
\usepackage[polish]{babel}
\usepackage[margin=2.2cm]{geometry}
\usepackage{enumitem}
\usepackage{booktabs,tabularx,array}

\newcolumntype{Y}{>{\raggedright\arraybackslash}X}
\newcommand{\IK}{\textsuperscript{\textbf{[IK]}}}
\newcommand{\HW}{\textsuperscript{\textbf{[HW]}}}
\newcommand{\WS}{\textsuperscript{\textbf{[wsp.]}}}

\setlength{\parindent}{0pt}
\setlength{\parskip}{4pt}
\setlist{itemsep=0pt,topsep=2pt}
\pagestyle{empty}

\begin{document}

\section*{Plan prac}

Termin zakończenia: koniec listopada 2026. Plan jest podzielony na etapy, a~nie
na sztywne zadania. Daty są orientacyjne i~będą aktualizowane po każdym spotkaniu.

\subsection*{Stan na start (październik)}
\begin{itemize}
  \item wstępny przegląd literatury \WS
  \item pierwsza wersja generatora danych syntetycznych \IK
  \item środowisko treningowe i~pierwsze modele, wstępne porównanie na danych syntetycznych \IK
  \item kompletowanie sprzętu i~koncepcja stanowiska \HW
\end{itemize}

\subsection*{Etapy}
\begin{center}
\small
\begin{tabularx}{\linewidth}{>{\raggedright\arraybackslash}p{3.0cm}YY}
\toprule
\textbf{Etap} & \textbf{Zadania} & \textbf{Rezultat} \\
\midrule
\textbf{1. Rozwój metody}\newline\footnotesize ok.\ 20.10--1.11 &
  \IK{} rozbudowa generatora, kolejne warianty modeli i~eksperymenty\newline
  \HW{} budowa stanowiska, oprogramowanie akwizycji\newline
  \WS{} szkielet pracy, rozdziały teoretyczne &
  działające stanowisko; wybrane kierunki rozwoju modelu \\
\addlinespace
\textbf{2. Dane rzeczywiste}\newline\footnotesize ok.\ 2--15.11 &
  \HW{} kalibracja, pomiary parametrów kamer, nagranie zbioru rzeczywistego\newline
  \IK{} dopasowanie generatora do zmierzonych parametrów, trening wybranych modeli\newline
  \WS{} opis generatora i~stanowiska &
  zbiór rzeczywisty; modele gotowe do ewaluacji \\
\addlinespace
\textbf{3. Ewaluacja}\newline\footnotesize ok.\ 16--22.11 &
  \WS{} testy na danych syntetycznych i~rzeczywistych, porównanie modeli, wnioski\newline
  \WS{} rozdziały z~wynikami &
  komplet wyników; wersja robocza pracy \\
\addlinespace
\textbf{4. Redakcja}\newline\footnotesize ok.\ 23--30.11 &
  \WS{} poprawki po uwagach opiekuna, uzupełnienia, finalizacja &
  praca gotowa do oddania \\
\bottomrule
\end{tabularx}
\end{center}

\subsection*{Zasady}
\begin{itemize}
  \item Spotkania kontrolne z~opiekunem po każdym etapie (lub częściej, w~razie potrzeby).
  \item Pisanie pracy przebiega równolegle z~pracami od pierwszego etapu.
  \item Najważniejsza zależność między częściami to parametry kamer zmierzone na stanowisku.
        Generator i~modele są przez nie dostrajane, więc pomiary mają priorytet w~etapie~2.
  \item Rezerwa czasowa jest w~etapach 3--4. Gdy coś się opóźni, ograniczamy liczbę wariantów
        modeli, a~nie ewaluację na danych rzeczywistych.
\end{itemize}

\end{document}
